% STATUS: DRAFT
\chapter{Anatomy of the result}\label{ch:anatomy}

\begin{physmotiv}
How does the abstract entropy $-\tau(\hat\rho\log\hat\rho)$ of the type II$_1$ algebra turn into the geometric formula $A/4G_N+S_{\rm out}$? The mechanism is the finite-dimensional identity of \cref{ch:entropy_II}, together with the de Sitter first law: matter and observer energy inside the patch reduce the horizon area. This chapter lays out which states are normal, what ``semiclassical'' means, why the observer's mass matters, and why the entropy is bounded above by that of empty de Sitter.
\end{physmotiv}

\section{Semiclassical states}

A \emph{semiclassical state} of $\A_{\rm obs}$ is a vector of the form
\begin{equation}
\hat\Phi=\Phi\otimes g^{1/2}(x),\qquad\Phi\in\HH\ \text{(a QFT state of finite energy relative to }\Psi_{\rm dS}),
\end{equation}
with $g$ a smooth, slowly varying probability density in $x=-\beta_{\rm dS}q$: slowly varying on the scale set by the modular Hamiltonian fluctuations of $\Phi$, so that $\Phi$ looks like a small perturbation of $\Psi_{\rm dS}$ compared with the spread of the observer's energy. These states are normal for $\A_{\rm obs}$. In the finite-dimensional model of \cref{ch:entropy_II} the entropy is exactly
\begin{equation}
S(\hat\Phi)=h(g)+\langle X\rangle_{\hat\Phi}-S_{\rm rel}(\Phi\|\Psi_{\rm dS}),\qquad h(g)=-\int g\log g,\quad X=K+x.
\end{equation}
For type III$_1$ $\A_R$ the same formula holds, to leading order in the semiclassical expansion, with Araki's relative entropy (\cref{ch:modham}); this is the content of \cite{CLPW2023}.

\section{From the formula to generalized entropy}

For a state excited only inside the patch ($\Phi$ differs from $\Psi_{\rm dS}$ by a unitary of $\A_R$) the relative entropy and the first law (\cref{ch:modham}) give
\begin{equation}
S_{\rm rel}(\Phi\|\Psi_{\rm dS})=\Delta\langle K_R\rangle_\Phi-\Delta S_{\rm out},\qquad K_R=\beta_{\rm dS}H_R,
\end{equation}
with $\Delta S_{\rm out}$ the change of the QFT entanglement entropy of the exterior (a finite difference) and $\langle K\rangle_\Phi=\Delta\langle K_R\rangle$ for such a state. Since $X=K+x$,
\begin{equation}
S(\hat\Phi)=h(g)+\langle x\rangle_g+\Delta S_{\rm out}.
\end{equation}
Now the gravitational input. The constraint $\hat H=H+q=0$ ties the observer's energy $q$ to the energy $E$ of the excitation inside the patch, and to first order a total energy $E$ inside the static patch lowers the horizon area by $\delta(A/4G_N)=-\beta_{\rm dS}E$ (cosmological first law, \cref{ch:ds_special}). The clock term $\langle x\rangle=-\beta_{\rm dS}\langle q\rangle$ is therefore, by the constraint, the area change of the horizon. Up to the entropy $h(g)$ of the clock's own distribution (subleading for broad $g$) and a state-independent constant,
\begin{equation}
S(\hat\Phi)\;\simeq\;\frac{A_{\rm dS}}{4G_N}+\delta\frac{A}{4G_N}+\Delta S_{\rm out}+\text{const},
\end{equation}
which is the generalized entropy $\Sgen=A/4G_N+S_{\rm out}$ of the patch in the state. The clock term $h(g)$ is the entropy of the observer's own energy distribution, which is subleading when $g$ is broad (semiclassical regime). This heuristic sketch is how the result is organized; the careful derivation, including the role of the observer's mass, the precise regime of validity and corrections, is in CLPW \cite{CLPW2023} and (for general subregions) in JSS \cite{JSS2023}.

\begin{keyidea}
Entropy of a semiclassical state $=$ \emph{clock term} (energy of the observer, which by the constraint equals minus the matter energy, hence the area change) $+$ \emph{bulk entropy difference}. The crossed product knows the area term because the constraint knows the first law.
\end{keyidea}

\section{Properties}

\begin{itemize}
\item \textbf{Upper bound.} $S(\hat\Phi)\le0$ for all states in $\A_{\rm dS}$, with equality only for $\Psi_{\max}$. Therefore (given the identification above) $\Sgen\le S_{\rm dS}$ for all semiclassical states: \emph{empty de Sitter space has the largest generalized entropy}, as expected geometrically since any matter in the patch lowers the horizon area.
\item \textbf{Dependence on the observer's mass.} The observer contributes its rest energy to $q$, hence lowers the horizon area (a Schwarzschild--de Sitter-like correction): a heavier observer has $A<A_{\rm dS}$. The distribution $g$ must be compatible with $q\ge0$; the maximum entropy state is the distribution $\beta\ee^{-\beta q}$ for the energy of the observer, i.e.\ an observer in thermal equilibrium with the horizon.
\item \textbf{Normal states.} States with finite $\langle X\rangle$, finite $S_{\rm rel}$ and $g$ supported at $q\ge0$ are in the folium of $\tau(P\,\cdot\,)$. States with sharp $q=q_0$ are not normal (not in $L^2$).
\item \textbf{$\beta$ dependence.} $\beta_{\rm dS}=2\pi\ell$ enters through $K=\beta H$ and $x=-\beta q$; changing $\ell$ rescales the unit of $q$ but not the algebra's structure.
\end{itemize}

\begin{whatwrong}
\begin{itemize}
\item The step ``clock energy $\Leftrightarrow$ area change'' is heuristic. In CLPW it is derived by solving the constraint for the semiclassical metric perturbation; the factors of $G_N$ and the identification of the zero of entropy are not fixed by the algebra.
\item The approximation $S\simeq h(g)+\langle x\rangle+\Delta S_{\rm out}$ drops terms suppressed by $1/\beta q$ etc.; the algebraic entropy is exact within the algebra, but the \emph{geometric interpretation} is semiclassical.
\item The relative-entropy term dominates only for states with energy $\ll$ observer energy fluctuations. For states with a large number of excitations the semiclassical expansion breaks down and the geometry must be backreacted.
\end{itemize}
\end{whatwrong}

\section*{Exercises}
\begin{exercise}
For $g(x)=\lambda\ee^{\lambda x}\theta(-x)$ show $S=1-1/\lambda-\log\lambda$ (\cref{ch:entropy_II}) and interpret $\lambda\ne1$ as the observer being at a temperature different from the horizon's.
\end{exercise}
\begin{exercise}
Using the first law $\delta(A/4G)=-\beta E$ and $S_{\rm rel}=\beta\Delta\langle H\rangle-\Delta S_{\rm out}\ge0$, show $\delta\Sgen=-S_{\rm rel}\le0$ for states with $\Phi$ obtained by acting with a unitary of the patch on $\Psi_{\rm dS}$. Interpret.
\end{exercise}
\begin{exercise}
Explain why the maximum-entropy state is not a state of the QFT alone but includes the observer's thermal distribution.
\end{exercise}
